\chapter{Chernobyl: Root-Cause Analysis}
\label{ch:analysis}

Why did Chernobyl happen, and what did the world conclude? The answer evolved
significantly between 1986 and 1992, and the evolution is itself instructive: it is
a case study in how the apportionment of blame between ``human error'' and
``design flaw'' depends on how deeply one understands the machine. This chapter
analyses the accident through the feedback theory of Parts~II and~III and traces
the shift from INSAG-1 to INSAG-7.

\section{The accident as a feedback failure}

In the language of this book, Chernobyl was the failure of a closed-loop control
system whose loop gain became positive. Recall the stability principle: a reactor
is stable when rising temperature/power removes reactivity faster than it is added.
At Chernobyl, in its low-power, low-ORM, distorted-flux state, three feedback
elements all had the wrong sign or were too weak:

\begin{itemize}[nosep]
\item The \textbf{void coefficient} was strongly positive: rising power boiled
water, which added reactivity (Chapters~\ref{ch:feedback},~\ref{ch:rbmk}).
\item The \textbf{prompt Doppler feedback}, though negative, was too weak at low
power to overcome the void term (the fuel-temperature swings at low power are
small).
\item The \textbf{scram}, intended as the ultimate negative reactivity, was
momentarily \emph{positive} because of the graphite displacers (the positive scram
effect).
\end{itemize}

With the dominant feedbacks positive, the Nyquist locus of Chapter~\ref{ch:linear}
enclosed the critical point, and the Nordheim--Fuchs balance of
Chapter~\ref{ch:nonlinear} had no benign turnover: there was no net negative
prompt feedback to cap the excursion. The accident is, in this sense, exactly what
the stability theory predicts for a machine with these coefficients driven into
this regime. It was not an inexplicable catastrophe; it was the equations
asserting themselves.

\section{Contributing causes: a layered failure}

No single fault caused the accident; it was a chain in which each link mattered.

\textbf{Design causes.} The positive void coefficient, the positive scram effect
of the graphite-tipped rods, the slow ($\sim 18\unit{s}$) shutdown, the large core
prone to spatial flux distortion, and the lack of automatic enforcement of the
ORM---all were properties of the machine.

\textbf{Operational causes.} The test was conducted at far too low a power; the
ORM was driven far below its minimum; protective trips were disabled; and the crew
proceeded despite the reactor being in an unfamiliar and unstable state.

\textbf{Organisational and cultural causes.} The positive scram effect was known
from the 1983 Ignalina event but neither corrected nor communicated; operating
documentation understated the danger of the low-ORM regime; the test procedure was
inadequate and inadequately reviewed; and a safety culture that treated such
margins as bureaucratic rather than physical allowed the boundary to be crossed.

The deepest lesson of Chernobyl is precisely that these layers combined. A machine
with the wrong sign of feedback can be operated safely \emph{only} if the
operators rigorously avoid the regimes where that sign bites---and that requires
the organisation to know, document, and enforce those boundaries. All three failed
together.

\section{From INSAG-1 to INSAG-7: the evolution of understanding}

The International Nuclear Safety Advisory Group (INSAG) issued two major reports.

\textbf{INSAG-1 (1986)}, drawing heavily on the initial Soviet presentation,
emphasised \emph{operator error}: the violations of procedure---running at low
power, defeating protections, depleting the ORM---were presented as the primary
cause. This framing placed the weight of responsibility on the crew.

\textbf{INSAG-7 (1992)}, after the Soviet authorities released fuller information
(including the State Committee's and the Soviet scientists' own deeper analyses),
substantially revised the picture. It emphasised \emph{design deficiencies}---the
positive void coefficient and especially the positive scram effect---as root
causes, while still acknowledging the operational violations. In INSAG-7's
assessment, the operators were drawn into a trap that the reactor's physics and the
inadequate information about it had set. Its often-quoted observation is that the
conditions under which the positive scram effect would be dangerous were believed
never to occur, yet ``appeared in almost every detail'' at Chernobyl.

\begin{remark}[Why the shift matters]
The movement from ``the operators broke the rules'' to ``the machine had flaws
that made rule-breaking catastrophic, and the rules themselves were poorly
justified to those who had to follow them'' is the central interpretive lesson of
Chernobyl. It reframes safety from a matter of operator discipline alone to a
matter of \emph{inherent design}: a reactor should not be one procedure violation
away from prompt criticality, and its shutdown system should never add reactivity.
Modern safety philosophy---inherent and passive safety, defence in depth, and a
strong safety culture---descends directly from this reframing.
\end{remark}

\section{The fixes to the RBMK fleet}

After the accident the surviving RBMK reactors were extensively modified to address
the very feedback flaws this book has emphasised:
\begin{itemize}[nosep]
\item Fuel enrichment was increased (from $\sim 2\%$ toward $2.4\%$) and additional
fixed absorbers were installed, which \emph{reduced the void coefficient} so that
the power coefficient became negative across the operating range.
\item The control rods were redesigned and the displacer geometry changed to
\emph{eliminate the positive scram effect}, with faster insertion and a larger
number of rods kept inserted.
\item A hard minimum operating reactivity margin was enforced, with automatic
protection, removing the low-ORM trap.
\item Improved instrumentation and faster emergency protection were added.
\end{itemize}
These changes did not make the RBMK a Western-style reactor, but they removed the
specific positive-feedback pathways that caused the accident---an implicit
confirmation that those pathways were indeed the root cause.

\section{Comparison with Western practice}

It is worth stating plainly the contrast that makes Chernobyl so instructive.
Western light-water reactors are designed with a \emph{negative} void coefficient
and a \emph{negative} power coefficient, achieved by under-moderation; their
shutdown systems insert only negative reactivity, quickly; and their containment
buildings are designed to retain radioactivity even given core damage---a feature
the RBMK lacked. The same reactivity-insertion event that destroyed Chernobyl would
be self-limited in a PWR by its prompt negative feedback (the BORAX/SPERT result),
and any residual release would be contained. Chernobyl is thus simultaneously a
story about one flawed design and a vindication of the inherent-safety principles
embodied in others.

\keypoint{Chernobyl was a feedback failure: in its operating state the dominant
reactivity feedbacks were positive, so the coupled system was unstable and the
prompt excursion had no benign limit. The understanding evolved from blaming
operators (INSAG-1) to recognising design flaws---the positive void coefficient
and positive scram effect---as root causes (INSAG-7). The enduring lesson is that
safety must be inherent in the physics, not contingent on perfect operation.}


\section*{Exercises}
\addcontentsline{toc}{section}{Exercises}
\begin{enumerate}[leftmargin=*,itemsep=2pt]
\item Frame the accident as a feedback-loop failure using the language of Chapters 6--9.
\item Summarise the differences between INSAG-1 and INSAG-7 and why the shift matters.
\item List the design, operational, and organisational causes and argue why all three were necessary.
\item Which post-accident RBMK modifications addressed the void coefficient, and how?
\item Contrast the RBMK with a Western PWR on void coefficient, scram, and containment.
\end{enumerate}
